\section{Abstract}

\textbf{DREAM-Net} (Dynamic Recall and Elastic Adaptive Memory) — это исследовательская архитектура нейронных сетей, предлагающая альтернативный подход к обработке последовательностей, основанный на принципах предиктивного кодирования и синаптической пластичности.

В отличие от распространённых архитектур, в которых параметры модели фиксируются после этапа обучения, DREAM-Net исследует возможность интеграции механизмов адаптации непосредственно в цикл инференса. Модель предполагает, что разделение весов на «быстрые» (адаптивные) и «медленные» (стабильные), в сочетании с иерархической передачей ошибки предсказания, может способствовать более эффективной обработке нестационарных данных, таких как речевые сигналы.

Архитектура включает три основных компонента:
\begin{enumerate}
    \item \textbf{Иерархическое предиктивное кодирование} — передача остаточной ошибки между слоями иерархии.
    \item \textbf{Двухмасштабная пластичность} — низкоранговое обновление быстрых весов ($U$) при сохранении стабильности медленных весов ($W$).
    \item \textbf{Статистический контроль пластичности} — модуляция скорости обучения на основе оценки новизны входного сигнала.
\end{enumerate}

Представленная спецификация описывает математическую формулировку архитектуры и служит основой для последующей экспериментальной проверки гипотез об её эффективности, устойчивости и вычислительной сложности.

\section{Мотивация и предпосылки исследования}

\subsection{Ограничения статических архитектур}

Современные подходы к обработке последовательностей, включая Трансформеры и модели пространства состояний (SSM), демонстрируют высокие результаты на задачах с фиксированным распределением данных. Однако ряд работ указывает на потенциальные ограничения таких архитектур в условиях нестационарности:

\begin{itemize}
    \item \textbf{Отсутствие онлайн-адаптации:} После обучения параметры модели не изменяются, что может требовать периодического дообучения при смене условий.
    \item \textbf{Вычислительная сложность контекста:} Механизмы внимания часто имеют квадратичную зависимость от длины последовательности.
    \item \textbf{Пластичность-стабильность дилемма:} Адаптация к новым данным может приводить к деградации ранее усвоенных знаний.
\end{itemize}

Биологические нейронные системы, в отличие от искусственных, демонстрируют способность к непрерывной адаптации без катастрофического забывания. Хотя прямое копирование биологических механизмов на кремниевые архитектуры невозможно, изучение их функциональных принципов может предложить новые направления для инженерных решений.

\subsection{Биологически вдохновлённые принципы как источник гипотез}

DREAM-Net не претендует на биологическую реалистичность. Вместо этого архитектура заимствует ряд функциональных идей, которые могут быть полезны с инженерной точки зрения:

\begin{table}[h]
\centering
\caption{Биологические принципы и их инженерная интерпретация}
\label{tab:bio_principles}
\begin{tabular}{p{3.5cm}p{4.5cm}p{4cm}}
\toprule
\textbf{Биологический принцип} & \textbf{Инженерная интерпретация в DREAM-Net} & \textbf{Исследовательский вопрос} \\
\midrule
Предиктивное кодирование & Иерархия, передающая только ошибку предсказания & Может ли такая схема снизить избыточность вычислений? \\
Синаптическая пластичность (STDP) & Низкоранговое обновление быстрых весов $U$ & Позволяет ли это ускорить адаптацию к локальным изменениям? \\
Нейромодуляция & Статистический гейт пластичности (Surprise Gate) & Может ли избирательная пластичность улучшить устойчивость к шуму? \\
Консолидация памяти & Асинхронный перенос знаний из $U$ в $W$ & Способствует ли это долгосрочному удержанию навыков? \\
\bottomrule
\end{tabular}
\end{table}

Каждый из этих компонентов формулируется как гипотеза, требующая экспериментальной проверки.

\subsection{Позиционирование исследования}

DREAM-Net предлагается как исследовательская платформа для изучения динамических архитектур, а не как готовая замена существующим решениям.

\textbf{Цели работы:}
\begin{enumerate}
    \item Предложить математически корректную формулировку архитектуры с онлайн-пластичностью.
    \item Исследовать устойчивость такой системы к нестационарным входным данным.
    \item Оценить вычислительную эффективность по сравнению с базовыми архитектурами (LSTM, SSM, Transformer) на контрольных задачах.
\end{enumerate}

\textbf{Ограничения:}
\begin{itemize}
    \item Архитектура не претендует на универсальность; её применимость может быть ограничена задачами с определённой структурой временных зависимостей.
    \item Все заявления об эффективности носят предварительный характер и требуют валидации на реальных датасетах.
    \item Биологические аналогии используются исключительно как эвристический ориентир, а не как доказательство корректности.
\end{itemize}

\section{Структура спецификации}

Данный документ описывает архитектуру DREAM-Net в формате технической спецификации. Последующие разделы включают:

\begin{itemize}
    \item \textbf{Mathematical Formulation:} Формальные уравнения предиктивного кодирования, пластичности и контроля удивления.
    \item \textbf{Control Logic:} Описание конечного автомата режимов работы и условий переходов.
    \item \textbf{Memory Architecture:} Детали реализации двухмасштабной памяти и механизма консолидации.
    \item \textbf{Safety \& Concurrency:} Меры обеспечения численной устойчивости и потокобезопасности.
    \item \textbf{Experimental Plan:} План валидации, включая метрики, базлайны и абляционные тесты.
\end{itemize}

\section{Заключение раздела}

Представленная спецификация фиксирует текущее состояние разработки архитектуры DREAM-Net. Все описанные механизмы являются предметом дальнейшего исследования. Экспериментальные результаты, если они будут получены, будут представлены отдельно с указанием условий воспроизведения и статистической значимости.
